\clearpage
\phantomsection
\addcontentsline{toc}{section}{参考文献}

\begingroup
    \zihao{5}\songti
    \setlength{\baselineskip}{20pt}
    
    \ctexset{
        section = {
            format = \raggedright\zihao{-3}\heiti\bfseries,
            beforeskip = 0pt,
            afterskip = 0.5\baselineskip,
            indent = 0pt,
            break = {},
        }
    }
    \renewcommand{\refname}{参考文献}
    
    \begin{thebibliography}{99}
        \setlength{\topsep}{0pt}
        \setlength{\partopsep}{0pt}
        \addtolength{\itemsep}{-0.5em}
        \bibitem{ref1} 刘国钧，陈绍业，王凤翥，等. 图书馆目录[M]. 北京：高等教育出版社，1957：15-18.
        \bibitem{ref2} 李晓东，张庆红，叶瑾琳. 气候学研究的若干理论问题[J]. 北京大学学报：自然科学版，1999, 35(1)：101-106.
        \bibitem{ref3} 辛希孟. 信息技术与信息服务国际研讨会论文集：A集[C]. 北京：中国社会科学出版社，1994.
        \bibitem{ref4} 钟文发. 非线性规划在可燃毒物配置中的应用[C]//赵玮. 运筹学的理论与应用：中国运筹学
        \bibitem{ref5} 会第五届大会论文集. 西安：西安电子科技大学出版社，1996：468-471.
        \bibitem{ref6} 张筑生. 微分半动力系统的不变集[D]. 北京：北京大学，1983.
        \bibitem{ref7} 冯西桥. 核反应堆压力管道与压力容器的 LBB 分析[R]. 北京：清华大学，1997.
        \bibitem{ref8} 谢希德. 创造学习的新思路[N]. 人民日报，1998-12-25(10).
        \bibitem{ref9} 中华人民共和国国家技术监督局. GB/T 16159—1996，汉语拼音正词法基本规则[S]. 北京：中国标准出版社，1996.
        \bibitem{ref10} 姜锡洲. 一种温热外敷药制备方案[P]. 中国：881056073，1989-07-26.
        \bibitem{ref11} 王明亮. 关于中国学术期刊标准化数据库系统工程的进展[EB/OL]. [1998-10-04]. http://www.cajcd.edu.cn/pub/wml.tex/980810-2.html.
    \end{thebibliography}
\endgroup

\vspace{1\baselineskip}
\begingroup
    \zihao{5}\songti\color{red}
    \noindent 【参考文献的序号左顶格，并用数字加括号表示，如 [1]，[2]，...，与正文中的指示序号格式一致。每一参考文献条目的最后均以“.”结束。字符、标点的全角、半角状态应尽量统一设置。参考文献条目超过一行的，所列文献折行应左缩进与第一行的首字(母)对齐。所引文献作者不多于 3 人时，作者应全部列出。合作者多于 3 人时，只列前 3 人，其后加“等”字或其他与之相应的字。】
\endgroup
